\subsection{Reconstruction from overlaps}

Here $\mathcal L_n(Y_m)$ denotes the set of length-$n$ blocks occurring in
$Y_m$.

\begin{lemma}[A signed Fibonacci gap shifts by one index]
\label{lem:signed-fibonacci-gap-shift}
Let $m\ge 2$ and let $d_1,\ldots,d_{m-1}\in\{-1,0,1\}$.  If
\[
\sum_{k=1}^{m-1} d_kF_{k+1}=F_m,
\]
then
\[
\sum_{k=1}^{m-1} d_kF_k=F_{m-1}.
\]
If instead
\[
\sum_{k=1}^{m-1} d_kF_{k+1}=-F_m,
\]
then
\[
\sum_{k=1}^{m-1} d_kF_k=-F_{m-1}.
\]
\end{lemma}

\begin{proof}
It is enough to prove the positive case, since the negative case follows by
replacing each $d_k$ with $-d_k$.

Let
\[
\varphi:=\frac{1+\sqrt5}{2},
\qquad
\psi:=\frac{1-\sqrt5}{2}.
\]
By Binet's formula,
\[
F_n=\frac{\varphi^n-\psi^n}{\sqrt5}
\qquad (n\ge 1).
\]
Define
\[
P(z):=\sum_{k=1}^{m-1} d_kz^k.
\]
The hypothesis reads
\[
\frac{\varphi P(\varphi)-\psi P(\psi)}{\sqrt5}=F_m
=
\frac{\varphi^m-\psi^m}{\sqrt5},
\]
so
\[
C:=\varphi P(\varphi)-\varphi^m
=
\psi P(\psi)-\psi^m.
\]
The first expression is an algebraic integer in $\ZZ[\varphi]$, while the
second is its Galois conjugate.  Since they are equal, $C$ is fixed by
conjugation and therefore lies in $\mathbb{Q}$; hence $C$ is an ordinary
integer.

On the other hand,
\[
|C|
\le
\sum_{k=1}^{m-1} |\psi|^{k+1}+|\psi|^m
<
\sum_{j=2}^{\infty} |\psi|^j
=
1.
\]
Thus the integer $C$ satisfies $|C|<1$, so $C=0$.  Therefore
\[
\varphi P(\varphi)=\varphi^m,
\qquad
\psi P(\psi)=\psi^m,
\]
and hence
\[
P(\varphi)=\varphi^{m-1},
\qquad
P(\psi)=\psi^{m-1}.
\]
Applying Binet's formula again gives
\[
\sum_{k=1}^{m-1} d_kF_k
=
\frac{P(\varphi)-P(\psi)}{\sqrt5}
=
\frac{\varphi^{m-1}-\psi^{m-1}}{\sqrt5}
=
F_{m-1}.
\]
\end{proof}

\begin{proposition}[Pairwise overlap reduction]
\label{prop:pair-overlap-reduction}
For every $m\ge 3$ there is a well-defined map
\[
\Theta_m:\mathcal L_2(Y_m)\to X_{m-1}
\]
with the following property.  If
\[
a=\Fold_m(pr),
\qquad
b=\Fold_m(rq)
\]
for some bits $p,q\in\{0,1\}$ and some overlap word
$r=r_1\cdots r_{m-1}\in\Omega_{m-1}$, then
\[
\Theta_m(ab)=\Fold_{m-1}(r).
\]
\end{proposition}

\begin{proof}
Take two raw words
\[
prq,\qquad p'r'q'
\]
of length $m+1$ such that
\[
\Fold_m(pr)=\Fold_m(p'r'),
\qquad
\Fold_m(rq)=\Fold_m(r'q').
\]
Write
\[
\Delta:=N(r)-N(r').
\]
Since $r,r'\in\Omega_{m-1}$, one has
\[
|\Delta|\le \sum_{k=1}^{m-1}F_{k+1}=F_{m+2}-2<F_{m+2}.
\]
By Corollary~\ref{cor:visible-value-overflow}, equality of the second folded
symbols implies
\[
\Delta+(q-q')F_{m+1}\in\{-F_{m+2},0,F_{m+2}\}.
\]
Because $q-q'\in\{-1,0,1\}$ and $F_{m+2}=F_{m+1}+F_m$, this forces
\[
\Delta\in\{0,\pm F_m,\pm F_{m+1}\}.
\]

We claim that the cases $\Delta=\pm F_m$ are impossible.  Suppose first that
\[
\Delta=F_m.
\]
Set
\[
d_k:=r_k-r_k'\in\{-1,0,1\}
\qquad (1\le k\le m-1).
\]
Then
\[
\sum_{k=1}^{m-1} d_kF_{k+1}=F_m,
\]
so Lemma~\ref{lem:signed-fibonacci-gap-shift} gives
\[
\sum_{k=1}^{m-1} d_kF_k=F_{m-1}.
\]
Hence
\[
\begin{aligned}
N(pr)-N(p'r')
&=
(p-p')+\sum_{k=1}^{m-1} d_kF_{k+2}\\
&=
(p-p')+\sum_{k=1}^{m-1} d_k(F_{k+1}+F_k)\\
&=
(p-p')+F_m+F_{m-1}\\
&=
(p-p')+F_{m+1}.
\end{aligned}
\]
This lies in
\[
\{F_{m+1}-1,F_{m+1},F_{m+1}+1\},
\]
which is disjoint from $\{-F_{m+2},0,F_{m+2}\}$ because
$F_{m+2}=F_{m+1}+F_m>F_{m+1}+1$.  By
Corollary~\ref{cor:visible-value-overflow}, the first folded symbols cannot be
equal, a contradiction.  The case $\Delta=-F_m$ is identical after changing
signs.

Therefore
\[
\Delta\in\{0,\pm F_{m+1}\}.
\]
By Corollary~\ref{cor:visible-value-overflow} at length $m-1$, this is exactly
the condition
\[
\Fold_{m-1}(r)=\Fold_{m-1}(r').
\]
Thus the folded overlap depends only on the pair of output symbols, and the map
$\Theta_m$ is well defined.
\end{proof}

\begin{remark}
The restriction $m\ge 3$ is essential.  For $m=2$, both raw words $00$ and
$11$ fold to $00$, so the pair $(00,00)\in\mathcal L_2(Y_2)$ does not
determine the overlapping bit.
\end{remark}

\begin{definition}[Recursive overlap decoders]
\label{def:recursive-overlap-decoder}
For integers $m\ge 3$ and $1\le n\le m$, define
\[
\Theta_{m,n}:\mathcal L_n(Y_m)\to X_{m-n+1}
\]
as follows.  For $n=1$, set
\[
\Theta_{m,1}(a_1):=a_1.
\]
For $(m,n)=(3,2)$, set
\[
\Theta_{3,2}(a_1a_2):=\Theta_3(a_1a_2).
\]
For $(m,n)=(3,3)$, let $\omega_1\cdots\omega_5$ be the unique raw lift from
Proposition~\ref{prop:G3-base-case} and define
\[
\Theta_{3,3}(a_1a_2a_3):=\omega_3.
\]
For $m\ge 4$ and $2\le n\le m$, define
\[
\Theta_{m,n}(a_1\cdots a_n)
:=
\Theta_{m-1,n-1}\bigl(
\Theta_m(a_1a_2)\,
\Theta_m(a_2a_3)\cdots
\Theta_m(a_{n-1}a_n)
\bigr).
\]
\end{definition}

\begin{proposition}[Common-overlap decoding]
\label{prop:common-overlap-decoder}
Let $m\ge 3$ and $1\le n\le m$.  Suppose
\[
\omega=\omega_1\cdots\omega_{m+n-1}\in\Omega_{m+n-1},
\qquad
a_j:=\Fold_m(\omega_j\cdots\omega_{j+m-1})
\quad (1\le j\le n).
\]
Then
\[
\Theta_{m,n}(a_1\cdots a_n)
=
\Fold_{m-n+1}(\omega_n\cdots\omega_m).
\]
Thus $\Theta_{m,n}$ recovers the fold of the common overlap of the $n$
consecutive raw windows.
\end{proposition}

\begin{proof}
We argue by induction on $m$.

If $m=3$, there are three cases.  For $n=1$, the claim is the definition of
$\Theta_{3,1}$.  For $n=2$, it is exactly
Proposition~\ref{prop:pair-overlap-reduction}.  For $n=3$,
Proposition~\ref{prop:G3-base-case} gives a unique raw lift
\[
\omega_1\cdots\omega_5
\]
of $a_1a_2a_3$, and by definition
\[
\Theta_{3,3}(a_1a_2a_3)=\omega_3.
\]
Since $X_1=\{0,1\}$ and $\Fold_1$ fixes admissible words by
Proposition~\ref{prop:fold-basic}\textup{(ii)}, this is
\[
\Fold_1(\omega_3).
\]
Thus the statement holds for $m=3$.

Now assume $m\ge 4$ and that the proposition is known at level $m-1$ for every
admissible block length.  The case $n=1$ is again the definition.  Assume
$n\ge 2$.  By Proposition~\ref{prop:pair-overlap-reduction}, for each
$1\le j\le n-1$ one has
\[
\Theta_m(a_ja_{j+1})
=
\Fold_{m-1}(\omega_{j+1}\cdots\omega_{j+m-1}).
\]
These are precisely the $(m-1)$-folds of the $n-1$ consecutive length-$(m-1)$
windows cut from the raw block
\[
\omega_2\cdots\omega_{m+n-2}.
\]
Applying the induction hypothesis at level $m-1$ to the $n-1$ folded
length-$(m-1)$ windows yields
\[
\Theta_{m,n}(a_1\cdots a_n)
=
\Theta_{m-1,n-1}\bigl(
\Theta_m(a_1a_2)\cdots\Theta_m(a_{n-1}a_n)
\bigr)
=
\Fold_{m-n+1}(\omega_n\cdots\omega_m),
\]
as claimed.
\end{proof}

\begin{theorem}[Unique lift of a length-$m$ output block]
\label{thm:length-m-block-unique-lift}
Let $m\ge 3$.  Every block
\[
a_1\cdots a_m\in\mathcal L_m(Y_m)
\]
has a unique raw lift
\[
\omega_1\cdots\omega_{2m-1}\in\Omega_{2m-1}
\]
satisfying
\[
\Fold_m(\omega_j\cdots\omega_{j+m-1})=a_j
\qquad (1\le j\le m).
\]
\end{theorem}

\begin{proof}
Existence is immediate from the definition of $\mathcal L_m(Y_m)$.

For uniqueness, suppose
\[
\omega_1\cdots\omega_{2m-1}
\]
is any lift of the given output block.  For $1\le k\le m$ define
\[
L_k:=\Theta_{m,m-k+1}(a_1\cdots a_{m-k+1}),
\qquad
R_k:=\Theta_{m,m-k+1}(a_k\cdots a_m).
\]
By Proposition~\ref{prop:common-overlap-decoder},
\[
L_k=\Fold_k(\omega_{m-k+1}\cdots\omega_m),
\qquad
R_k=\Fold_k(\omega_m\cdots\omega_{m+k-1}).
\]

For $k=1$ this gives
\[
L_1=R_1=\omega_m\in\{0,1\},
\]
so the central bit is determined.

Now fix $2\le k\le m$ and assume inductively that the central blocks
\[
\omega_{m-k+2}\cdots\omega_m,
\qquad
\omega_m\cdots\omega_{m+k-2}
\]
have already been determined.  The word $L_k$ is the fold of the raw $k$-block
\[
\omega_{m-k+1}\omega_{m-k+2}\cdots\omega_m,
\]
whose suffix of length $k-1$ is known.  By
Proposition~\ref{prop:one-step-left-resolving}, applied with $k$ in place of
$m$, the two possible prepended bits $0$ and $1$ produce distinct folds.
Therefore $L_k$ determines the new left bit $\omega_{m-k+1}$ uniquely.

Similarly, $R_k$ is the fold of
\[
\omega_m\cdots\omega_{m+k-2}\omega_{m+k-1},
\]
whose prefix of length $k-1$ is known.  By
Proposition~\ref{prop:one-step-right-resolving}, again with $k$ in place of
$m$, the appended bit is uniquely determined by $R_k$.  Hence $\omega_{m+k-1}$
is fixed as well.

Proceeding outward from $k=1$ to $k=m$ reconstructs all bits
\[
\omega_1,\ldots,\omega_{2m-1}
\]
uniquely.
\end{proof}

\begin{corollary}[Uniform synchronizing length]
\label{cor:uniform-synchronizing-length}
Let $m\ge 3$.  Every block in $\mathcal L_m(Y_m)$ is synchronizing for the
presentation $G_m$.  Moreover the block
\[
e_{m,0}e_{m,1}\cdots e_{m,m-2}
\]
of length $m-1$ is not synchronizing.  Consequently the minimal synchronizing
length of $G_m$ is exactly $m$.
\end{corollary}

\begin{proof}
By Theorem~\ref{thm:length-m-block-unique-lift}, every block in
$\mathcal L_m(Y_m)$ is carried by exactly one length-$m$ path in $G_m$, hence
is synchronizing.

For the second claim, Proposition~\ref{prop:one-sided-inverse-memory-lower-bound}
constructs two sequences $U_m$ and $V_m$ whose first $m-1$ output symbols are
\[
e_{m,0}e_{m,1}\cdots e_{m,m-2}
\]
and coincide there.  The corresponding length-$(m-1)$ paths in $G_m$ end at the
suffix vertices
\[
010^{m-3}
\qquad\text{and}\qquad
10^{m-2},
\]
which are distinct.  Thus the block admits at least two lifts with different
terminal vertices, so it is not synchronizing.  The distinct $m$th output
symbols constructed in Proposition~\ref{prop:one-sided-inverse-memory-lower-bound}
show in particular that these two lifts do not merge after one further step.
\end{proof}

\begin{theorem}[Finite-block losslessness in Fibonacci numeration]
\label{thm:block-bijection-threshold}
Let $m\ge 3$ and $n\ge m$.  Define
\[
B_{m,n}:\Omega_{n+m-1}\to\mathcal L_n(Y_m),
\qquad
B_{m,n}(\omega_1\cdots\omega_{n+m-1}):=a_1\cdots a_n,
\]
where
\[
a_j:=\Fold_m(\omega_j\cdots\omega_{j+m-1})
\qquad (1\le j\le n).
\]
Then $B_{m,n}$ is a bijection.
\end{theorem}

\begin{proof}
Surjectivity is immediate from the definition of $\mathcal L_n(Y_m)$.

For injectivity, fix a block
\[
a_1\cdots a_n\in\mathcal L_n(Y_m).
\]
Its prefix
\[
a_1\cdots a_m\in\mathcal L_m(Y_m)
\]
has a unique raw lift by Theorem~\ref{thm:length-m-block-unique-lift}, hence it
is carried by a unique length-$m$ path in $G_m$.  Since $G_m$ is right-resolving
by Proposition~\ref{prop:one-step-right-resolving}, each subsequent symbol
$a_{m+1},\ldots,a_n$ determines the next edge uniquely.  Thus the whole block
$a_1\cdots a_n$ is carried by a unique length-$n$ path in $G_m$.

Now a length-$n$ path in $G_m$ is uniquely determined by its initial vertex
$v=v_1\cdots v_{m-1}\in V_m$ together with the appended bits
$b_1,\ldots,b_n\in\{0,1\}$.  Equivalently, it is uniquely determined by the raw
word
\[
v_1\cdots v_{m-1}b_1\cdots b_n\in\Omega_{n+m-1}.
\]
By construction of the labels in $G_m$, this raw word is precisely the unique
length-$(n+m-1)$ block whose consecutive folded windows are
$a_1,\ldots,a_n$.  Therefore $B_{m,n}$ is injective.
\end{proof}

\begin{corollary}[Exact block counts above the threshold]
\label{cor:block-counts-threshold}
For every $m\ge 3$ and $n\ge m$,
\[
\lvert\mathcal L_n(Y_m)\rvert=2^{n+m-1}.
\]
In particular,
\[
\lvert\mathcal L_m(Y_m)\rvert=2^{2m-1}.
\]
\end{corollary}

\begin{proof}
Theorem~\ref{thm:block-bijection-threshold} identifies $\mathcal L_n(Y_m)$
with $\Omega_{n+m-1}=\{0,1\}^{n+m-1}$, whose cardinality is $2^{n+m-1}$.
\end{proof}

\begin{table}[t]
\centering
\small
\begin{tabular}{ccc}
\toprule
$(m,n)$ & direct enumeration of $\lvert\mathcal L_n(Y_m)\rvert$ & $2^{n+m-1}$ \\
\midrule
$(3,3)$ & $32$ & $32$ \\
$(3,4)$ & $64$ & $64$ \\
$(3,5)$ & $128$ & $128$ \\
$(4,4)$ & $128$ & $128$ \\
$(4,5)$ & $256$ & $256$ \\
$(4,6)$ & $512$ & $512$ \\
\bottomrule
\end{tabular}
\caption{Small-scale verification of Corollary~\ref{cor:block-counts-threshold}
obtained by brute-force enumeration of all raw blocks in $\Omega_{n+m-1}$ and
folding each consecutive length-$m$ window.}
\label{tab:small-block-counts}
\end{table}

\begin{corollary}[Higher-block identification with the full two-shift]
\label{cor:higher-block-identification}
Let $m\ge 3$ and $n\ge m$.  The bijection $B_{m,n}$ identifies the
$n$-block presentation of $Y_m$ with the $(n+m-1)$-block presentation of the
full binary shift.  Equivalently, consecutive blocks in $\mathcal L_n(Y_m)$
correspond exactly to consecutive overlapping blocks in $\Omega_{n+m-1}$.
\end{corollary}

\begin{proof}
Let
\[
a=a_1\cdots a_n,
\qquad
a'=a_2\cdots a_{n+1}
\]
be consecutive $n$-blocks in some point of $Y_m$.  If
\[
B_{m,n}^{-1}(a)=\omega_1\cdots\omega_{n+m-1},
\qquad
B_{m,n}^{-1}(a')=\omega'_1\cdots\omega'_{n+m-1},
\]
then both raw words are lifts of the same length-$(n+1)$ output block
$a_1\cdots a_{n+1}\in\mathcal L_{n+1}(Y_m)$.  By
Theorem~\ref{thm:block-bijection-threshold}, that length-$(n+1)$ output block
has a unique raw lift in $\Omega_{n+m}$, so necessarily
\[
\omega'_1\cdots\omega'_{n+m-2}=\omega_2\cdots\omega_{n+m-1}.
\]
Thus adjacency of $n$-blocks in $Y_m$ is exactly overlap by $n+m-2$ bits of
their raw lifts in $\Omega_{n+m-1}$, which is the adjacency rule in the
$(n+m-1)$-block presentation of the full binary shift.  Hence the two
higher-block presentations are canonically identified.
\end{proof}

\begin{corollary}[Linear-time finite-block decoding]
\label{cor:linear-time-block-decoding}
Fix $m\ge 3$.  For every $n\ge m$, a block
\[
a_1\cdots a_n\in\mathcal L_n(Y_m)
\]
can be decoded to its unique raw lift in deterministic time $O_m(n)$ and
working space $O(m)$ beyond the space used to store the output.
\end{corollary}

\begin{proof}
By Theorem~\ref{thm:length-m-block-unique-lift}, the initial block
$a_1\cdots a_m$ determines the unique raw lift
\[
\omega_1\cdots\omega_{2m-1}.
\]
This fixes the current suffix vertex
\[
\omega_{m+1}\cdots\omega_{2m-1}\in V_m.
\]
For each $j=m+1,\ldots,n$, the current suffix vertex together with the next
symbol $a_j$ determines a unique outgoing edge of $G_m$ by
Proposition~\ref{prop:one-step-right-resolving}; following that edge reveals
the next appended bit $\omega_{j+m-1}$.  Repeating this step decodes the entire
raw lift.  Since $m$ is fixed, the initialization cost is $O_m(1)$, each new
output symbol is processed in constant time, and only the current
length-$(m-1)$ suffix together with the already decoded raw word must be
stored.
\end{proof}

\begin{remark}[Small decoding examples]
\label{rem:small-decoding-examples}
For $m=3$, the output block
\[
(a_1,a_2,a_3)=(001,010,100)
\]
has unique raw lift $00100$.  Thus the central-bit inverse returns
\[
\kappa_{3,1}(001,010,100)=1.
\]
For $m=4$, the output block
\[
(a_1,a_2,a_3,a_4)=(1001,0010,1000,0000)
\]
has unique raw lift $1110000$.  Hence the position-parameterized inverse family
reads off
\[
\kappa_{4,0}=1,\qquad
\kappa_{4,1}=1,\qquad
\kappa_{4,2}=1,\qquad
\kappa_{4,3}=0
\]
on this single normalized block.  The first example is already fully
admissible, while the second passes through nontrivial local normalization
before decoding.
\end{remark}

\begin{theorem}[Exact finite-type order over the folded alphabet]
\label{thm:Y-m-exact-SFT-order}
For every $m\ge 3$, the shift $Y_m$ is an $m$-step shift of finite type over
the alphabet $X_m$, and it is not an $(m-1)$-step shift of finite type.
\end{theorem}

\begin{proof}
Let
\[
y=(y_t)_{t\in\ZZ}\in X_m^{\ZZ}
\]
be such that every length-$(m+1)$ block of $y$ lies in
$\mathcal L_{m+1}(Y_m)$.  For each $t\in\ZZ$, the length-$m$ block
\[
b_t:=y_{t-m+1}\cdots y_t
\]
belongs to $\mathcal L_m(Y_m)$ and is therefore synchronizing by
Corollary~\ref{cor:uniform-synchronizing-length}.  Let $v_t$ denote the unique
terminal vertex of a path in $G_m$ labeled by $b_t$.

Now fix $t$.  Since
\[
y_{t-m+1}\cdots y_ty_{t+1}\in\mathcal L_{m+1}(Y_m),
\]
there exists a path in $G_m$ carrying this length-$(m+1)$ block.  Its prefix
of length $m$ is labeled by $b_t$, so by uniqueness its terminal vertex is
$v_t$.  Its suffix of length $m$ is labeled by
\[
b_{t+1}=y_{t-m+2}\cdots y_{t+1},
\]
so by the same uniqueness statement its terminal vertex is $v_{t+1}$.  Hence
the last edge of that path is an edge
\[
v_t\xrightarrow{\,y_{t+1}\,}v_{t+1}.
\]
Thus the vertices $(v_t)_{t\in\ZZ}$ form a bi-infinite path in $G_m$ labeled by
$y$, and therefore $y\in Y_m$.  This proves that $Y_m$ is $m$-step.

To show that $Y_m$ is not $(m-1)$-step, let
\[
b:=e_{m,0}e_{m,1}\cdots e_{m,m-2}\in\mathcal L_{m-1}(Y_m).
\]
By Corollary~\ref{cor:uniform-synchronizing-length}, the block $b$ is not
synchronizing, so reading $b$ in $G_m$ can end at at least two distinct
vertices.  By Proposition~\ref{prop:one-sided-inverse-memory-lower-bound}, the
length-$m$ block
\[
bc=b\,010^{m-2}
\]
belongs to $\mathcal L_m(Y_m)$.  By
Theorem~\ref{thm:length-m-block-unique-lift}, this length-$m$ block is carried
by exactly one path in $G_m$.  Consequently, among the terminal vertices
reached by reading $b$, exactly one admits an outgoing edge labeled
$010^{m-2}$.  Since $b$ is not synchronizing, at least one other terminal
vertex reached by $b$ does not admit that label.

Every finite path in $G_m$ occurs in some bi-infinite labeled path because the
underlying graph is the binary de Bruijn graph on words of length $m-1$ and is
therefore strongly connected.  Hence there exist two occurrences of the same
length-$(m-1)$ block $b$ in points of $Y_m$ such that one occurrence may be
followed by $010^{m-2}$ and another may not.  Thus the admissible next symbol
after $b$ depends on more than the last $m-1$ output symbols, so $Y_m$ is not
an $(m-1)$-step shift of finite type.
\end{proof}

\begin{theorem}[Local inverse rules]
\label{thm:Phi-m-conjugacy-threshold}
For every $m\ge 3$ and every $0\le r\le m-1$, let
\[
\kappa_{m,r}:\mathcal L_m(Y_m)\to\{0,1\}
\]
send a block $a_1\cdots a_m$ to the $(r+1)$st bit of its unique raw lift from
Theorem~\ref{thm:length-m-block-unique-lift}.  Define
\[
\Psi_m^{(r)}:Y_m\to\{0,1\}^{\ZZ},
\qquad
(\Psi_m^{(r)}(y))_t:=\kappa_{m,r}(y_{t-r}\cdots y_{t-r+m-1}).
\]
Then $\Psi_m^{(r)}$ is a sliding block code with memory $r$ and anticipation
$m-1-r$, and
\[
\Psi_m^{(r)}\circ\Phi_m=\mathrm{id}_{\{0,1\}^{\ZZ}},
\qquad
\Phi_m\circ\Psi_m^{(r)}=\mathrm{id}_{Y_m}.
\]
Equivalently, each length-$m$ observation window yields an exact inverse rule
for $\Phi_m:\{0,1\}^{\ZZ}\to Y_m$, and $\Phi_m$ is a topological conjugacy
between the full binary shift and the image shift $Y_m$.
\end{theorem}

\begin{proof}
By definition, $(\Psi_m^{(r)}(y))_t$ depends only on the block
$y_{t-r}\cdots y_{t-r+m-1}$, so $\Psi_m^{(r)}$ has memory $r$ and anticipation
$m-1-r$.

Let $s\in\{0,1\}^{\ZZ}$ and set $y:=\Phi_m(s)$.  For each $t\in\ZZ$, applying
Theorem~\ref{thm:length-m-block-unique-lift} to the output block
\[
y_{t-r}\cdots y_{t-r+m-1}
\]
shows that its unique raw lift is
\[
s_{t-r}\cdots s_{t-r+2m-2}.
\]
Therefore the $(r+1)$st bit of that lift is $s_t$, so
\[
(\Psi_m^{(r)}(y))_t=s_t.
\]
Hence
\[
\Psi_m^{(r)}(\Phi_m(s))=s,
\]
so $\Psi_m^{(r)}\circ\Phi_m=\mathrm{id}$.

Now let $y\in Y_m$.  By definition of $Y_m$ there exists
$s\in\{0,1\}^{\ZZ}$ with $\Phi_m(s)=y$.  Applying the first identity to this
$s$ gives
\[
\Psi_m^{(r)}(y)=\Psi_m^{(r)}(\Phi_m(s))=s,
\]
and therefore
\[
\Phi_m(\Psi_m^{(r)}(y))=\Phi_m(s)=y.
\]
Thus $\Phi_m$ and $\Psi_m^{(r)}$ are mutual inverses.
\end{proof}

\begin{corollary}[The one-sided memory threshold is exact]
\label{cor:Phi-m-memory-threshold-exact}
Let $m\ge 3$.  Among all sliding block codes
\[
\Psi:Y_m\to\{0,1\}^{\ZZ}
\]
with anticipation $0$ and
\[
\Psi\circ\Phi_m=\mathrm{id}_{\{0,1\}^{\ZZ}},
\]
the smallest possible memory is exactly $m-1$.
\end{corollary}

\begin{proof}
The endpoint inverse $\Psi_m^{(m-1)}$ from
Theorem~\ref{thm:Phi-m-conjugacy-threshold} has memory $m-1$, while
Proposition~\ref{prop:one-sided-inverse-memory-lower-bound} shows that no
one-sided left inverse can use smaller memory.
\end{proof}
